% ----------  Chapter 9 : Contingency-Planning Policies ----------
\chapter{Contingency-Planning Policies}
\label{chap:contingency}

This chapter defines a coherent set of policies that strengthen
GreenGrid Energy’s capability to prepare for, respond to, and
recover from disruptive events.  These policies operationalise the risk-management
commitments adopted by the Board in Chapters 3–7.
%------------------------------------------------------------------
\section{Policy Framework}
\label{sec:framework}

\begin{longtable}{@{}p{0.28\linewidth}p{0.45\linewidth}p{0.22\linewidth}@{}}
\toprule
\textbf{Policy} & \textbf{Purpose} & \textbf{ISO/IEC 27001:2022 Annex A}\\
\midrule
\endhead
\ac{GCPP} &
Enterprise-level principles governing \emph{all} contingency capabilities. &
Foundational (applies to the whole standard)\\[2pt]

\ac{IRPP} &
Preparation and coordination for information-security incidents. &
5.24 \textit{Information-security incident management planning \& preparation}\\[2pt]

\ac{DRPP} &
Restoration of ICT/OT services to meet defined RTO/RPO. &
5.29 \textit{Information-security continuity}\\[2pt]

\ac{BCPP} &
Continuity of critical business processes beyond ICT recovery. &
5.30 \textit{ICT readiness for business continuity}\\
\bottomrule
\end{longtable}

The policies are \emph{cumulative}: the GCPP sets universal rules, while
\ac{IRPP}, \ac{DRPP} and \ac{BCPP} add discipline-specific requirements.

%------------------------------------------------------------------
\section{\ac{GCPP}}
\subsection*{Purpose}
Safeguard life, environment, assets and reputation by
\emph{anticipating, preventing and managing} disruptive events.

\subsection*{Scope}
All personnel, sites, OT and IT systems (including SCADA, ERP and cloud
environments) inside the ISMS boundary defined in Chapter~2.

\subsection*{Policy Statements}
\begin{enumerate}[label=\arabic*.]
  \item \textbf{Risk-based planning}\,—All contingency plans shall derive
        from documented risk assessments and the Statement of Applicability.
  \item \textbf{Single source of truth}\,—An approved master copy of each
        plan shall reside in the secure document repository with version
        control and change-management records.
  \item \textbf{Hierarchy \& integration}\,—Department-level plans (IR, DR,
        BC) must align with and reference the \ac{GCPP}.
  \item \textbf{Testing cadence}\,—Each plan shall be exercised at least
        annually and after material changes or significant incidents.
  \item \textbf{Continuous improvement}\,—Lessons learned from exercises or
        live events feed the corrective-action register and trigger updates.
\end{enumerate}

\subsection*{Roles \& Responsibilities}

\begin{longtable}{@{}p{0.25\linewidth}p{0.7\linewidth}@{}}
\toprule
\textbf{Role} & \textbf{Key duties}\\
\midrule
\endhead
Board of Directors & Approve the \ac{GCPP} and resource its implementation.\\[2pt]
\ac{CISO} & Own the framework and coordinate
audits.\\[2pt]
Department Directors & Maintain department-specific plans and ensure staff
readiness.\\[2pt]
All Employees \& Contractors & Participate in training, report incidents,
follow approved procedures.\\
\bottomrule
\end{longtable}

\subsection*{Training \& Awareness}
Mandatory induction plus refresher training every 12 months; scenario
workshops for Engineering \& Grid Operations owing to \ac{SCADA} criticality.

\subsection*{Review \& Maintenance}
The CISO schedules semi-annual management-review meetings to confirm plan
adequacy and regulatory alignment.

%------------------------------------------------------------------
\section{\ac{IRPP}}
\label{sec:IRPP}
\subsection*{Purpose}
Formalise preparation, escalation and coordination required to manage
information-security incidents rapidly and effectively.

\subsection*{Scope}
Real or suspected compromises of confidentiality, integrity or availability
affecting IT or OT (e.g.\ \ac{SCADA} manipulation, ERP data breach, phishing).

\subsection*{Policy Requirements}
\begin{enumerate}[label=\arabic*.]
  \item \textbf{\ac{IRT}} consisting of: Incident
        Commander, SOC Analysts, OT-Security Engineer, Communications
        Officer and Legal Adviser.
  \item \textbf{Seven-phase method}\,—Preparation $\rightarrow$ Detection
        $\rightarrow$ Assessment $\rightarrow$ Containment
        $\rightarrow$ Eradication $\rightarrow$ Recovery
        $\rightarrow$ Post-incident review.
  \item \textbf{15-minute reporting rule}\,—Employees must inform the SOC or
        on-call \ac{IRT} within 15 minutes of discovery.
  \item \textbf{Classification \& escalation}\,—Severity 1 events (critical
        infrastructure impact) are escalated to the CEO and regulators
        within one hour.
  \item \textbf{Forensic evidence protection}\,—Maintain chain-of-custody
        logs for potential legal proceedings.
  \item \textbf{Metrics}\,—MTTD, MTTC and number of repeat incidents.
\end{enumerate}

\subsection*{Testing}
Table-top exercise every six months and at least one live
\emph{purple-team} drill per year.

%------------------------------------------------------------------
\section{\ac{DRPP}}
\label{sec:DRPP}
\subsection*{Purpose}
Restore essential ICT and OT services—including SCADA control, cloud
analytics and ERP—to their agreed \ac{RTO} and
\ac{RPO} after a disruption.

\subsection*{Scope}
All production data centres, public-cloud regions and OT substations that
support core revenue streams.

\subsection*{Policy Requirements}
\begin{enumerate}[label=\arabic*.]
  \item \textbf{Business Impact Analysis}\,—Review critical processes
        annually to set or revise RTO/RPO.
  \item \textbf{Redundant architecture}\,—Dual-site or multi-region design
        for SCADA masters and ERP; hot-standby OT–IT gateways.
  \item \textbf{Backup strategy}\,—3-2-1 rule with immutable object
        storage.
  \item \textbf{Invocation authority}\,—Only the COO or CISO may declare a
        disaster and activate the DR plan.
  \item \textbf{Test schedule}\,—Functional fail-over test (quarterly);
        full blackout recovery simulation (annually).
  \item \textbf{Supplier coordination}\,—DR provisions embedded in SLA
        clauses with critical vendors.
\end{enumerate}

%------------------------------------------------------------------
\section{\ac{BCPP}}
\label{sec:BCPP}
\subsection*{Purpose}
Ensure that \emph{critical business processes}—energy distribution,
customer billing, supply-chain procurement, R\&D labs—continue at predefined
capacity levels during and after disruptive events.

\subsection*{Scope}
All functions above plus supporting HR, Legal and Administration services
necessary for compliance and workforce welfare.

\subsection*{Policy Requirements}
\begin{enumerate}[label=\arabic*.]
  \item \textbf{Continuity-management lifecycle}: Context \& BIA
        $\rightarrow$ Strategy $\rightarrow$ Plan Development
        $\rightarrow$ Training $\rightarrow$ Exercising
        $\rightarrow$ Maintenance.
  \item \textbf{Minimum Service Levels (MSL)}: \\
        \begin{tabular}{@{}l l@{}}
        Grid-Operations Control & $\le$ 15 minutes downtime; 80\,\% load
        capacity.\\
        Customer Contact Centre & 75 \% calls handled within 4 hours.\\
        Procurement Portal & Emergency ordering within 8 hours.\\
        \end{tabular}
  \item \textbf{Alternative facilities}: Remote SCADA console at secondary
        NOC; cloud-native CRM \emph{warm} tenant for customer support.
  \item \textbf{People continuity}: Cross-training matrix for key roles to
        mitigate single-point-of-knowledge risk identified in
        Chapter~6.
  \item \textbf{Communication plan}: Pre-approved templates for regulators,
        media and customers, vetted by Legal \& Administration.
\end{enumerate}

\subsection*{Exercises \& Evaluation}
Departmental walk-throughs every quarter; enterprise-wide simulation
(including key suppliers) every 24 months.  Corrective actions are tracked
in the Risk-Treatment Register.

%------------------------------------------------------------------
\section{Governance, Monitoring \& Review}
\begin{itemize}[leftmargin=*]
  \item The \emph{Risk \& Compliance Committee} receives quarterly KPI
        dashboards for IR, DR and BC performance.
  \item Internal audit evaluates each policy’s effectiveness against
        ISO/IEC 27001 clauses and Annex A controls.
  \item Policies are reviewed at least annually or whenever major business,
        technology or regulatory changes occur.
\end{itemize}

%------------------------------------------------------------------
\section{Document-Control Table}
\label{sec:doccontrol}

\begin{longtable}{@{}p{0.28\linewidth}p{0.12\linewidth}p{0.07\linewidth}%
                  p{0.15\linewidth}p{0.15\linewidth}@{}}
\toprule
\textbf{Document} & \textbf{Owner} & \textbf{Ver.} &
\textbf{Approval Date} & \textbf{Next Review}\\
\midrule
\endhead
GCPP & CISO & 1.0 & 18-May-2025 & 18-May-2026\\
IRPP & CISO & 1.0 & 18-May-2025 & 18-May-2026\\
DRPP & COO/CISO & 1.0 & 18-May-2025 & 18-May-2026\\
BCPP & COO/CISO & 1.0 & 18-May-2025 & 18-May-2026\\
\bottomrule
\end{longtable}

% ----------  End of Chapter 9  ----------
